% Kapitel 2 --- Hintergrund und verwandte Arbeiten
% Budget: 2,5 Seiten

\section{Hintergrund und verwandte Arbeiten}
\label{sec:hintergrund}

\subsection{Informationsqualität als Rahmen}

Informationsqualität ist kein einzelnes Merkmal, sondern ein Bündel von Dimensionen.
\textcite{wang1996beyond} leiten aus Befragungen fünfzehn Dimensionen ab und ordnen sie in
intrinsische, kontextuelle, repräsentationale und zugangsbezogene Kategorien. Der
Standard ISO/IEC 25012 \parencite{iso25012} fasst dieselbe Idee normativ und definiert
unter anderem \emph{Consistency} als Freiheit von Widerspruch und Kohärenz mit anderen
Daten sowie \emph{Currentness} als richtiges Alter im Nutzungskontext. ISO/IEC 25024
\parencite{iso25024} ergänzt Messgrößen zu diesen Merkmalen.

Das Projekt untersucht fünf dieser Dimensionen getrennt: Konsistenz, Aktualität,
Auffindbarkeit, Verständlichkeit und Mehrsprachigkeit. Die vorliegende Arbeit bearbeitet die
Konsistenz, aus drei Gründen. Sie ist normativ verankert und braucht keine eigene Definition
des Qualitätsbegriffs; sie hat eine natürliche Analyseeinheit, den Sachverhalt, der sich
zählen lässt; und sie ist quellenübergreifend, lässt sich also nicht durch die Prüfung eines
einzelnen Dokuments erledigen. Wichtig ist dabei die Arbeitsteilung: Der Standard sagt,
\emph{was} gemessen wird, nicht \emph{wie}. Die Operationalisierung --- Extraktion,
Normalisierung, Typologie, Goldstandard --- ist der Beitrag dieser Arbeit und nicht aus der
Norm ableitbar.

\subsection{Widersprüche automatisch erkennen}

Der Referenzrahmen für automatische Widerspruchsprüfung ist die Aufgabenstellung von FEVER
\parencite{thorne2018fever}: Zu einer Behauptung wird eine Belegstelle gesucht und daraus ein
Urteil abgeleitet --- gestützt, widerlegt oder nicht entscheidbar. Der Datensatz umfasst
185.445 Behauptungen, die gegen Absätze aus Wikipedia geprüft werden. Für diese Arbeit ist
nicht der Umfang entscheidend, sondern die Struktur: Das Urteil hängt an einer zitierbaren
Stelle und nicht an der Überzeugung des Prüfers. Genau diese Struktur braucht die
Konsistenzmessung, denn sie vergleicht nicht eine Aussage mit Weltwissen, sondern zwei
Dokumente miteinander.

Für Fließtext, in dem Werte nicht exakt übereinstimmen, ist Natural Language Inference ein
Baustein. SNLI stellt rund 570.000 Satzpaare mit den Urteilen \emph{Implikation}, \emph{neutrale
Aussage} und \emph{Widerspruch} bereit \parencite{bowman2015snli}; XNLI erweitert das
Verfahren auf fünfzehn Sprachen \parencite{conneau2018xnli}. Der Nutzen für diese Arbeit ist
begrenzt und klar umrissen: NLI entscheidet, ob zwei Sätze einander widersprechen, wenn sie
verschieden formuliert sind. Die Prüfung, ob sie denselben Sachverhalt betreffen, ersetzt es
nicht.

Die Grenzen sind empirisch belegt. \textcite{pielka2022} zeigen, dass Widerspruchsmodelle an
Präpositionen und Verben, an Antonymen und Homonymen, an unvollständigen Sätzen, an langen
Absätzen und an seltenen Wörtern scheitern, und dass sich crowdsourced-maschinell übersetzte
Daten sprachlich von Textpaaren aus Internetquellen unterscheiden. Daraus folgt für diese
Arbeit zweierlei: Der Goldstandard muss aus dem eigenen Korpus stammen, und ein
SNLI-trainiertes Urteil ist kein Beleg für einen Widerspruch zwischen zwei JLU-Seiten. Es ist
ein Hinweis, der an der Fundstelle geprüft werden muss.

\subsection{Veraltung und Content Drift}

Zwei Phänomene werden häufig vermischt. \emph{Link Rot} bezeichnet Verweise, die ins Leere
führen; \emph{Content Drift} bezeichnet Verweise, die weiter funktionieren, aber auf
veränderten Inhalt zeigen \parencite{klein2014,jones2016}. Für diese Arbeit ist der zweite
Fall der relevante: Eine Seite bleibt erreichbar und sieht gepflegt aus, während die Angabe
darauf nicht mehr stimmt.

Wie ungleich solche Änderungen verteilt sind, zeigt \textcite{fetterly2003} an 150.836.209
HTML-Seiten über elf Wochen: Die Änderungsrate variiert stark zwischen Top-Level-Domains, und
größere Seiten ändern sich häufiger als kleinere. \textcite{cho2000freshness} zeigen
anschließend, wie sich Abrufpläne an die gemessene Änderungsrate anpassen lassen. Übertragen
auf Studieninformationen heißt das: Nicht jede Angabe veraltet gleich schnell, und ein
Vergleich über alle Angaben hinweg mischt schnelle und langsame Typen. Ob Fristen tatsächlich
häufiger abweichen als Modulnummern, ist eine der Hypothesen dieser Arbeit
(Abschnitt~\ref{sec:frage}); die Literatur zeigt nur, dass solche Unterschiede zu erwarten
sind.

Für die Fragestellung folgt daraus eine Unterscheidung, die die ganze Auswertung trägt: Ein
Teil der gefundenen Abweichungen ist kein Widerspruch, sondern eine nicht mitgezogene
Aktualisierung. Beide Fälle sehen im Datensatz gleich aus --- zwei verschiedene Werte für
denselben Sachverhalt --- und müssen dennoch verschieden behandelt werden. Die Typologie in
Abschnitt~\ref{sec:methode} trennt sie.

\subsection{Abgrenzung}

Drei Abgrenzungen sind wichtig. Verweise und Fristen \emph{innerhalb} der Ordnungen --- ob
eine Ordnung auf eine andere verweist und dabei die richtige Fassung nennt --- prüft die
Arbeit 2.2 desselben Projekts; hier geht es ausschließlich um Abweichungen \emph{zwischen}
Dokumenten. Ob zwei Erwähnungen dieselbe Entität meinen, ist Gegenstand von 2.4 (Entity
Linking); diese Arbeit setzt die Entitätsauflösung über eine gepflegte Aliasliste voraus,
statt sie zu bewerten. Und die Veraltung als eigene Dimension mit Zeitreihe und Vorhersage
ist 3.2 --- hier ist sie nur ein Erklärungsmuster für einen Teil der Abweichungen.
